%% Generated by scripts/build_manuscript.py from docs/paper/SUPPLEMENTARY.md -- do not edit.
\documentclass[11pt]{article}
\usepackage[margin=2.5cm]{geometry}
\usepackage{graphicx,booktabs,tabularx,amsmath,amssymb,xcolor,url}
\usepackage[authoryear,round]{natbib}
\renewcommand{\tabularxcolumn}[1]{>{\raggedright\arraybackslash}p{#1}}
\newcommand{\botrule}{\bottomrule}
\providecommand{\bibcommenthead}{}
\renewcommand{\thesection}{S\arabic{section}}
\renewcommand{\thetable}{S\arabic{table}}
\renewcommand{\thefigure}{S\arabic{figure}}
\title{Supplementary Information\\[4pt]\large The Price of Utility: Scientific Computation as Browser Admission Work, Measured Under Attack}
\date{}
\begin{document}
\maketitle

Additional file 1 for ``The Price of Utility: Scientific Computation as Browser Admission Work, Measured Under Attack''. Section and table numbers prefixed S refer to this file; all others refer to the main text.

\section{Experimental methods}

\textbf{Input preparation and qualification.} The five targets were fixed from the Vina-scoring results in Eberhardt et al.'s supplementary Table S2~\citep{Eberhardt21} before local docking. For each target, SHA-256 ordering with salt \texttt{published-vina-validation-2026-09-15:} selects 32 active and 64 decoy source compound IDs from DUD-E. All supplied states of each selected compound are retained; the crystal compound is excluded if its ID matches. There is no replacement after preparation or docking failure. Earlier DYR, AKT1 and PPARG failures remain negative results outside this qualified panel; this study does not establish that arbitrary targets pass stock Vina's gate.

Receptor waters and existing hydrogens are removed; REDUCE 3.16 with \texttt{-BUILD} rebuilds hydrogens and allows its standard flips. ADFRsuite's \texttt{prepare\_receptor4.py} uses \texttt{-U nphs\_lps\_waters}, retaining nonprotein chains; recognized metal charges are set to +2 in this protocol. Original DUD-E MOL2 ligand states pass through \texttt{prepare\_ligand4.py} defaults. Heavy-atom counts and coordinates are checked, and source/prepared files and tool hashes are retained. The box is 22 $\times$ 22 $\times$ 22 \AA{}, centered on the arithmetic centroid of the crystal ligand's heavy atoms, with 0.375 \AA{} grid spacing. This centroid convention and the small panel are explicit local choices. Stock Vina 1.2.7 is used rather than the publication's 1.2.0; this is not an exact reproduction of its full-library enrichment.

Stock qualification uses exhaustiveness 32, one CPU per job, nine modes and seed 104729. All 96 compounds must complete, with ROC-AUC at least 0.75, its bootstrap lower 95\% bound above 0.60, and EF10 at least 1.5. Redocking must have top-pose symmetry-corrected RMSD at most 2 \AA{} in at least two of seeds 104729, 130363 and 155921, without alignment. Only qualified targets enter the matched comparison. Screening score is the minimum over all required states; any missing state makes the compound incomplete.

\textbf{Compute matching and uncertainty.} For each ligand state and parent seed, the same-build monolithic E32 arm supplies its measured evaluation count E. The distributed arm runs \texttt{ceil(E / 256000)} independently seeded units, then uses the original finaliser. The same inputs, box and state aggregation are used in both arms. The matching variable is energy evaluations, not elapsed browser time; search-time and evaluation ratios are reported separately. This is an oracle-budget experimental comparison: production scheduling need not rerun stock docking to assign every budget. The 2,073 pairs represent state/seed jobs, not 2,073 independent compounds.

Paired ROC-AUC intervals use 5,000 stratified compound bootstrap resamples with random seed 104729. Active and decoy compound indices are sampled separately; the same sampled indices are applied to both arms and all three seeds before averaging seed-specific AUC differences. Repeated states and seeds are therefore not treated as independent compounds. AUC gives ties half credit. EF10 selects the top \texttt{ceil(0.1 $\times$ n)} compounds and gives fractional occupancy to ties at the cutoff. The predeclared non-inferiority margin is $-$0.05 AUC. Unit-level binomial intervals are descriptive and do not account for dependence among units from the same molecular job.

\textbf{Poisoning projection.} Measured score changes from the 33 re-finalised corpus jobs are pooled across targets, separately for each attack fraction. For each of 1,000 resampling draws (random seed 20260923), a change is sampled independently for each matched state/seed score. Compound minima and seed-averaged target AUCs are recomputed. This assumes exchangeable shifts across targets, states and activity labels; it does not model their correlation or an adversary selecting compounds by label. The resulting AUC changes are exploratory projections, not whole-panel attack measurements or confirmatory confidence intervals for real-world poisoning.

The underlying integrity experiment directly replaces or alters units in saved pools and reruns the finaliser; it does not exercise an admission-bypass path. At 5\% fabricated-energy units, 22 of 33 jobs worsen, by at most 1.859 kcal/mol; the maximum worsening of 2.632 occurs at 10\%. Across these pool-tampering experiments, the largest score improvement is 0.466 kcal/mol at 5\% fabricated-energy units. Recomputed output energies do not establish biological validity or exclude screening false positives.

\textbf{Availability timing and replication.} Queue operations use the prototype SQL, while arrivals, replay duration and puzzle costs are simulated. Honest arrivals are 0.5/s for 240 s, cycling three device classes (40 each), followed by 150 s of draining. Bootstrap attacks continue for all 390 s. Replay cost is fixed at 1.52 s. Honest puzzle times are exponential for a single puzzle or gamma-distributed for subpuzzles at measured device hash rates; attacker costs use expected hashes at 1.25 million hashes/s/core. Issuance is capped at 30 attacker attempts per 250 ms tick. Ten seconds is the honest bidding budget, not a hard end-to-end access deadline; random solve times can exceed it. Successful access within the observation window includes draining.

The original simulator processed replay completions only on 250 ms ticks, so a 1.52 s replay occupied a worker for 1.75 s. Amendment 12 processes completion and worker refill at exact service event times and dispatches available work after each input tick. Input and puzzle-completion polling remain at 250 ms, with the final input tick at 389.75 s if the simulation has not already drained. All 224 original grid cells are repeated at five fixed seeds, 20260925--20260929 (1,120 runs). Original results remain preserved. Private in-memory SQLite uses the same SQL and transaction boundaries to accelerate simulation; it is not a production database benchmark. Mean and min/max across seeds describe simulation variability, not uncertainty in measured rates or attacker optimality. Main-text Section 5.5 and Section S3 distinguish historical results from this corrected replication.

\section{First phone timing study}

\textbf{Device protocol.} A test page served from the benchmark's content delivery network runs real WebAssembly Vina units and hash puzzles on the same device: one cold unit, four warm calibration units, a measured JavaScript hash rate, then twelve rounds alternating units with puzzles so that thermal effects fall on both. Puzzle difficulty is calibrated on the device to the median warm unit. The page holds a screen wake lock and timestamps every hidden period; measurements overlapping a hidden period are flagged and excluded by rule. The server re-checks every unit's pool and trace hashes against the reference before storing a report. An earlier iPhone run on a previous page, which could not flag hidden periods, lost one unit to the Low Power Mode auto-lock (11.1 s against a 1.9 s median); it is reported here but excluded from all pooled figures. The budget phone's units rose from 7.4 to 10.1 s mid-run and recovered, consistent with thermal throttling; its first contribution took about 13 s (3.8 s asset delivery, 1.5 s restore, 7.5 s first unit).

Against a hashcash puzzle calibrated natively to the same median time, a docking unit had half the coefficient of variation (0.48 against 0.95) and a p99 of 3.38 s against 6.50 s. On four phones, a test page alternated real units with puzzles calibrated on the device to its median unit.

\begin{table}[htbp]
\caption{Docking unit and device-calibrated single hashcash puzzle on four phones, 12 rounds each.}
\centering
\small
\begin{tabularx}{\textwidth}{XXXX}
\toprule
Device & Unit median / max & Puzzle median / p95 / max & Units exact \\
\midrule
iPhone 15, Low Power Mode & 2.04 / 2.10 s & 1.85 / 6.57 / 8.24 s & 17/17 \\
Samsung Galaxy A57 & 2.02 / 2.10 s & 2.09 / 7.82 / 10.5 s & 17/17 \\
Moto Edge 50, Low Power Mode & 2.44 / 2.73 s & 1.86 / 8.83 / 10.4 s & 17/17 \\
Samsung Galaxy M30s (2019 budget) & 7.55 / 10.1 s & 5.86 / 36.1 / 53.2 s & 17/17 \\
\botrule
\end{tabularx}
\end{table}

Scaled to each device's median, 21.9\% of puzzle solves exceeded twice the median and none of 64 warm executions did. All 68 executions were bitwise identical to reference outputs across iOS and Android. They repeat four distinct reference units. The second study in main-text Section 5.2 likewise repeats four distinct units, producing 85 matching executions on five phones; execution counts should not be interpreted as counts of distinct workloads.

\section{Availability: prediction checks and replicated results}

Every prediction check from amendments 9, 9b and 10, original single-seed and corrected five-seed (amendment 12), with misses retained. The protocol labels the amendment-10 predictions A1--A6; here and in the main text they are T1--T6 so they cannot be confused with attacker behaviours A1--A8.

The ticket prototype underlying these simulations was also exercised through loopback HTTP with real puzzles and protocol operations, and a separate two-bundle native smoke check with actual full-unit replay accepted honest output and rejected corrupted trace commitments.

\begin{table}[htbp]
\caption{Original single-seed and corrected five-seed prediction checks. A check is a device-group outcome in an eligible grid cell; counts retain all misses.}
\centering
\small
\begin{tabularx}{\textwidth}{Xrr}
\toprule
Prediction & Original passed / eligible & Corrected passed / eligible \\
\midrule
C1, follow published price & 33/48 & 153/240 \\
C1, full-patience bid & 46/48 & 230/240 \\
T1, attested service below capacity & 120/120 & 600/600 \\
T2, saturation upper bounds & 112/120 & 557/600 \\
T3, anonymous capacity model & 118/118 & 590/590 \\
T4, inert attested lane & 24/24 & 120/120 \\
T5, trust bootstrap & 63/66 & 315/330 \\
T6, puzzle-free without fallback & 192/192 & 960/960 \\
\botrule
\end{tabularx}
\end{table}

\begin{table}[htbp]
\caption{Corrected eight-worker service, mean [minimum, maximum] across five seeds, 40 arrivals per device class per seed. These are observed ranges, not confidence intervals.}
\centering
\small
\begin{tabularx}{\textwidth}{XXXX}
\toprule
Mechanism, attacker cores & Budget & Mid-range & Flagship \\
\midrule
Fixed 18-bit, 0.25 & 1.000 [1.000, 1.000] & 1.000 [1.000, 1.000] & 1.000 [1.000, 1.000] \\
Fixed 18-bit, 1 & 0.945 [0.900, 0.975] & 0.940 [0.875, 1.000] & 0.915 [0.800, 0.975] \\
Fixed 18-bit, 4 & 0.390 [0.350, 0.425] & 0.325 [0.175, 0.475] & 0.365 [0.325, 0.400] \\
Fixed 18-bit, 16 & 0.230 [0.175, 0.275] & 0.145 [0.075, 0.250] & 0.195 [0.125, 0.250] \\
Full-patience priority, 0.25 & 0.995 [0.975, 1.000] & 0.965 [0.950, 0.975] & 1.000 [1.000, 1.000] \\
Full-patience priority, 1 & 0.990 [0.975, 1.000] & 0.970 [0.950, 0.975] & 0.965 [0.950, 0.975] \\
Full-patience priority, 4 & 0.070 [0.050, 0.100] & 0.965 [0.950, 0.975] & 1.000 [1.000, 1.000] \\
Full-patience priority, 16 & 0.110 [0.100, 0.125] & 0.135 [0.050, 0.200] & 0.130 [0.100, 0.175] \\
\botrule
\end{tabularx}
\end{table}

For direct historical comparison at four attacker cores, the original fixed-18-bit service fractions were 0.15 / 0.25 / 0.42 and priority fractions were 0.05 / 0.95 / 1.00. At sixteen cores these were 0.07 / 0.17 / 0.12 and 0.05 / 0.03 / 0.05, respectively. Original single-seed outcomes and new five-seed means differ in both timing and replication; the difference cannot be attributed solely to the service-time correction.

\begin{table}[htbp]
\caption{Trust bootstrap under sixteen attacker cores. Corrected entries are mean [minimum, maximum] across five seeds among token holders at 90\% coverage (36 per device class per seed).}
\centering
\small
\begin{tabularx}{\textwidth}{XXXX}
\toprule
Workers, attacker tokens/s & Budget & Mid-range & Flagship \\
\midrule
8, 0 & 1.000 [1.000, 1.000] & 1.000 [1.000, 1.000] & 1.000 [1.000, 1.000] \\
8, 0.1 & 1.000 [1.000, 1.000] & 1.000 [1.000, 1.000] & 1.000 [1.000, 1.000] \\
8, 1 & 1.000 [1.000, 1.000] & 1.000 [1.000, 1.000] & 1.000 [1.000, 1.000] \\
8, 10 & 0.639 [0.500, 0.750] & 0.944 [0.917, 1.000] & 0.861 [0.806, 0.889] \\
4, 10 & 0.311 [0.250, 0.389] & 0.189 [0.139, 0.250] & 0.100 [0.083, 0.139] \\
\botrule
\end{tabularx}
\end{table}

The original one-seed ten-token/s results were 0.92 / 0.89 / 0.97 at eight workers and 0.08 / 0.44 / 0.42 at four. The corrected results show substantial changes and seed variation, particularly for budget phones; increased replay capacity does not imply monotonic improvement for every group in this finite queue, fallback and retry process. These observations do not isolate which transient interaction causes each change.

\section{Further admission-economics configurations}

The committed scheduler was driven under a simulated clock with audit verdicts resampled from real replay records: 118 configurations, of which 110 ran to completion and 8 were stopped at 5.9 hours with their measured counts recorded. Values are attacker discount factors before the fixes unless stated.

\begin{table}[htbp]
\caption{Further admission-economics configurations on the committed scheduler.}
\centering
\small
\begin{tabularx}{\textwidth}{XXX}
\toprule
Configuration & Predicted & Measured \\
\midrule
Bundle tier, attacker always submits instead of walking away & as walking away & 0.42 \\
Attempt cap 1 / cap 16, one unit computed & 1.0 / 0.25 & 0.98 / 0.25 \\
Two audit draws, 2 or 3 units computed & 1.19 / 0.86 & 1.19 / 0.85 \\
Reduced-budget units, pool-only commitment, 64k / 128k evaluations & 0.51 / 0.60 & 0.54 / 0.60 \\
Same after reseeding & about 1.25 / 1.10 & 1.39 / 1.12 \\
Reduced-budget units, pool and trace commitment & never admitted & never admitted \\
Trusted tier, trust after 1 bundle, p = 0.1, immediate / deferred audit & about 0.5 / 0.61 & 0.60 / 0.63 \\
Disappearing workers holding 4 / 8 / 12 / 16 of 16 leases & refusals about D/16 & 0\% / 0\% / 0\% / 81\% refused \\
\botrule
\end{tabularx}
\end{table}

Trust acquisition. With the three-bundle threshold, acquiring trust by cached reuse cost 5--15 times honest work. At ten bundles, that strategy earned trust about 720 times per configuration from roughly 7,450 bundle admissions and never produced a trusted admission, because abandonments left every identity above the trusted tier's risk limit after the strategy's fixed 30-minute backoff; this bounds only a 30-minute-backoff attacker. On the enforced threshold at p = 0.1, an attacker planning trust after one bundle obtained no trusted admission across 3,000 identities while spending 12,000 units of honest work; one earning three bundles paid a factor of 2.04.

\bibliographystyle{sn-basic}
\bibliography{references}
\end{document}
